\section{Related Work}
\label{sec:related}

\textbf{Flight delay prediction.} Machine-learning work on flight delays spans classification and
regression, tree ensembles and deep models, and many prediction horizons and delay
thresholds~\cite{carvalho2021relevance}. Rebollo and Balakrishnan~\cite{rebollo2014characterization}
predict delays 2--24 hours ahead with random forests, using network-wide and link-level delay
states. Gui \emph{et al.}~\cite{gui2020flight} combine ADS-B trajectories, weather, schedules and
airport information, and report 90.2\% accuracy for binary classification with a random forest.
Choi \emph{et al.}~\cite{choi2016prediction} predict weather-induced delays using weather
forecasts. Yi \emph{et al.}~\cite{yi2021flight} stack several classifiers for a single airport and
handle class imbalance by oversampling. These studies draw on real-time, weather or trajectory
information, and they largely report accuracy or ROC-based metrics without assessing whether
predicted probabilities are calibrated.

\textbf{Delay propagation.} Delays spread through airport networks: Pyrgiotis
\emph{et al.}~\cite{pyrgiotis2013modelling} model how local congestion propagates and pushes demand
into the evening. This motivates our time-of-day, scheduled-congestion and past-only historical
delay features.

\textbf{Methods.} We use gradient-boosted trees~\cite{chen2016xgboost} alongside Logistic
Regression. Classifier scores are often poorly calibrated, and Platt scaling and isotonic
regression are standard corrections~\cite{niculescu2005predicting,guo2017calibration}. We report
the expected calibration error and reliability diagrams~\cite{guo2017calibration} and the Brier
score~\cite{brier1950verification}. Because delayed flights are the minority class, precision--recall
analysis is more informative than ROC analysis~\cite{saito2015precision}, so PR-AUC is our primary
metric.

\textbf{Leakage and time-aware evaluation.} Leakage, meaning the use of information that would not
be available when predicting, inflates offline results~\cite{kaufman2012leakage}. With temporally
dependent data, how a model is validated also matters~\cite{bergmeir2012use}. We use a chronological
hold-out because the model is deployed one day ahead and delay rates shift from month to month.

\textbf{Gap.} We study the stricter day-ahead setting with only the published schedule and
past-only outcomes, enforce leakage controls with tests, use each time period for a single
decision, and report calibrated probabilities alongside PR-AUC. This quantifies what such a model
can and cannot deliver.
